\documentclass[12pt,a4paper]{article}
\usepackage[utf8]{inputenc}
\usepackage[margin=1in]{geometry}
\usepackage{graphicx}
\usepackage{amsmath}
\usepackage{algorithm}
\usepackage{algorithmic}
\usepackage{hyperref}
\usepackage{cite}
\usepackage{listings}
\usepackage{xcolor}

\title{\textbf{Adaptive Packet Routing using Q-Learning in OMNeT++\\
\large A Reinforcement Learning Approach}}
\author{Q-Routing Implementation Project}
\date{\today}

\begin{document}

\begin{titlepage}
    \centering
    
    {\large CSE 4120: Technical Writing \& Seminar}
    \vspace{0.5cm}
    \begin{center}
    {\Large \bfseries \textsc{Adaptive Packet Routing using Q-Learning in OMNeT++: A Reinforcement Learning Approach}}
    \end{center}

    
    \vspace{3.5cm}
    \textbf{Submitted By}\\
    \vspace{.5cm}
    {\Large\bfseries Efty Hasan}\\
    {Roll: 2007052}
    
    \vspace{2.5cm}

    \textbf{Submitted To}\\
    \vspace{.5cm}
    
    {\Large\bfseries  Dr. K. M. Azharul Hasan}\\
    { Professor\\Department of Computer Science and Engineering\\Khulna University of Engineering \& Technology}
    
    \vspace{1cm}
    
    {\Large\bfseries  Kazi Saeed Alam}\\
    { Assistant Professor\\Department of Computer Science and Engineering\\Khulna University of Engineering \& Technology}
    
    \vfill 

    \includegraphics[width=0.15\textwidth]{kuet.png}\par
    
    
    {
    Department of Computer Science and Engineering\\
    % \vspace{0.3cm}
    Khulna University of Engineering \& Technology\\\
    % \vspace{0.3cm}
    Khulna 9203, Bangladesh\\
    }
    {\normalsize October 29, 2025\par}

\end{titlepage}

\maketitle

\section{Objective}

The primary objectives of this project are:

\begin{enumerate}
    \item \textbf{Implement Q-Routing Algorithm}: Develop a fully functional Q-Routing protocol in OMNeT++ that enables routers to learn optimal routing decisions through reinforcement learning.
    
    \item \textbf{Demonstrate Adaptive Learning}: Validate that routers can autonomously discover the fastest paths in a network with heterogeneous link delays without prior knowledge of network topology.
    
    \item \textbf{Verify Convergence}: Demonstrate that Q-values converge to reflect true end-to-end delays and that the system successfully exploits learned optimal paths.
    
    \item \textbf{Analyze Performance}: Evaluate the learning dynamics, exploration-exploitation trade-off, and overall system performance through simulation experiments.
    
    \item \textbf{Educational Implementation}: Create a clear, well-documented implementation that serves as a practical demonstration of reinforcement learning principles applied to network routing.
\end{enumerate}

\section{Introduction}

\subsection{Background}

Traditional routing protocols such as OSPF (Open Shortest Path First) and RIP (Routing Information Protocol) rely on static metrics and periodic updates to determine routing paths. These protocols face challenges in dynamic network environments where link conditions, traffic loads, and network topology change frequently. The need for adaptive routing mechanisms that can respond to real-time network conditions has motivated research into learning-based approaches.

\subsection{Q-Routing: A Reinforcement Learning Approach}

Q-Routing, introduced by Boyan and Littman in 1994 \cite{boyan1994packet}, applies reinforcement learning principles to the routing problem. Instead of relying on predefined metrics or global topology information, each router learns optimal forwarding decisions through experience, similar to how an agent learns optimal actions in a Markov Decision Process (MDP).

\subsubsection{Fundamental Concepts}

In Q-Routing, each router maintains a \textit{Q-table} that estimates the time required to deliver a packet to each destination through each neighboring router. The Q-value $Q_x(d,y)$ at router $x$ represents the estimated delivery time to destination $d$ via neighbor $y$.

\textbf{Key components:}
\begin{itemize}
    \item \textbf{State}: Current router and packet destination
    \item \textbf{Action}: Selection of next-hop neighbor
    \item \textbf{Reward}: Negative of delivery time (faster = better)
    \item \textbf{Policy}: Epsilon-greedy selection (exploitation vs. exploration)
\end{itemize}

\subsubsection{The Q-Learning Equation}

The core of Q-Routing is the Q-value update equation. When router $x$ forwards a packet destined for $d$ to neighbor $y$, and subsequently receives feedback about the delivery time, it updates its Q-value as:

\begin{equation}
Q_x(d,y) \leftarrow Q_x(d,y) + \alpha \left[ t + \min_{z} Q_y(d,z) - Q_x(d,y) \right]
\label{eq:qrouting}
\end{equation}

where:
\begin{itemize}
    \item $Q_x(d,y)$: Current Q-value estimate at router $x$ for destination $d$ via neighbor $y$
    \item $\alpha$: Learning rate (typically 0.5), controlling the weight of new information
    \item $t$: Transmission time from router $x$ to router $y$
    \item $\min_{z} Q_y(d,z)$: Minimum Q-value at next router $y$ for destination $d$
\end{itemize}

This equation implements the Bellman optimality principle: the optimal path cost equals the immediate cost plus the optimal cost from the next state.

\subsubsection{Simplified Q-Routing for Direct Connections}

In topologies where routers have direct connections to destinations (single-hop scenarios), Equation \ref{eq:qrouting} simplifies to:

\begin{equation}
Q_x(d,y) \leftarrow Q_x(d,y) + \alpha \left[ \text{RTT} - Q_x(d,y) \right]
\label{eq:simple}
\end{equation}

where RTT is the measured round-trip time. This simplification is valid because $\min_{z} Q_y(d,z) = 0$ at the destination, making the update purely based on observed delay.

\subsection{The Q-Routing Algorithm}

Algorithm \ref{alg:qrouting} presents the complete Q-Routing procedure:

\begin{algorithm}
\caption{Q-Routing Protocol}
\label{alg:qrouting}
\begin{algorithmic}[1]
\STATE \textbf{Initialize:} $Q_x(d,y) \leftarrow Q_{\text{init}}$ for all destinations $d$, neighbors $y$
\STATE \textbf{Initialize:} Explored gates set $E_x(d) \leftarrow \emptyset$ for all destinations $d$
\STATE
\WHILE{simulation running}
    \STATE \textbf{On receiving packet} destined for $d$:
    \STATE
    \IF{$\exists y \notin E_x(d)$} \textcolor{blue}{// Initial exploration phase}
        \STATE Select first unexplored neighbor $y$
        \STATE $E_x(d) \leftarrow E_x(d) \cup \{y\}$
    \ELSIF{$\text{random}() < \epsilon$} \textcolor{blue}{// Random exploration}
        \STATE Select random neighbor $y$
    \ELSE \textcolor{blue}{// Exploitation}
        \STATE $y \leftarrow \arg\min_{z} Q_x(d,z)$ \textcolor{blue}{// Choose best neighbor}
    \ENDIF
    \STATE
    \STATE Record: $\text{SendTime}_{\text{packet}} \leftarrow \text{CurrentTime}$
    \STATE Forward packet to neighbor $y$
    \STATE
    \STATE \textbf{On receiving response:}
    \STATE $\text{RTT} \leftarrow \text{CurrentTime} - \text{SendTime}_{\text{packet}}$
    \STATE $Q_x(d,y) \leftarrow Q_x(d,y) + \alpha \times (\text{RTT} - Q_x(d,y))$ \textcolor{blue}{// Update Q-value}
\ENDWHILE
\end{algorithmic}
\end{algorithm}

\subsection{Smart Exploration Strategy}

Our implementation incorporates a two-phase exploration strategy:

\begin{enumerate}
    \item \textbf{Initial Systematic Exploration}: Each possible action (neighbor) is tried at least once before exploitation begins. This ensures all paths are evaluated.
    
    \item \textbf{Epsilon-Greedy Exploration}: After initial exploration, the system uses epsilon-greedy policy with $\epsilon = 0.1$, meaning 10\% random exploration and 90\% exploitation of learned knowledge.
\end{enumerate}

This approach balances the exploration-exploitation dilemma, ensuring continued adaptation while primarily using the best known paths.

\section{Implementation Approach}

\subsection{Simulation Environment}

The implementation was developed using:
\begin{itemize}
    \item \textbf{OMNeT++ 6.2.0}: Discrete event simulation framework
    \item \textbf{Language}: C++ for module implementation, NED for network description
    \item \textbf{No INET Framework}: Pure OMNeT++ implementation for educational clarity
\end{itemize}

\subsection{Network Topology}

We designed a multi-router topology specifically to evaluate Q-Routing's path learning capabilities:

\begin{figure}[h]
\centering
\includegraphics[width=0.8\textwidth]{network_topology.png}
\caption{Network topology showing Source PC connected to three routers with different path delays}
\label{fig:network_topology}
\end{figure}

\begin{verbatim}
Source PC (addr=1)
    |---(Link: 1ms)---> Router1 (addr=11) ---(FastPath: 10ms)---> 
    |---(Link: 1ms)---> Router2 (addr=12) ---(MediumPath: 50ms)--->
    |---(Link: 1ms)---> Router3 (addr=13) ---(SlowPath: 100ms)---> 
                                        Destination PC (addr=2)
\end{verbatim}

This topology provides three parallel paths with different delays:
\begin{itemize}
    \item \textbf{Path 1 (via Router1)}: Total RTT = 22ms (1ms + 10ms + 10ms + 1ms)
    \item \textbf{Path 2 (via Router2)}: Total RTT = 102ms (1ms + 50ms + 50ms + 1ms)
    \item \textbf{Path 3 (via Router3)}: Total RTT = 202ms (1ms + 100ms + 100ms + 1ms)
\end{itemize}

\subsection{Module Design}

\subsubsection{Router Module (router.cc)}

The Router module implements the core Q-Routing algorithm:

\textbf{Key Data Structures:}
\begin{itemize}
    \item \texttt{map<long, map<int, double>> qTable}: Stores Q-values indexed by destination and gate
    \item \texttt{map<long, set<int>> exploredGates}: Tracks which gates have been tried for each destination
    \item \texttt{map<long, PacketInfo> pendingPackets}: Tracks sent packets for RTT calculation
\end{itemize}

\textbf{Key Methods:}
\begin{itemize}
    \item \texttt{selectGate(dest, inGate)}: Implements epsilon-greedy selection with smart exploration
    \item \texttt{updateQ(dest, gate, rtt, nextHopQ)}: Updates Q-values using Equation \ref{eq:simple}
    \item \texttt{handleMessage()}: Processes incoming packets, forwards data, learns from responses
\end{itemize}

\subsubsection{PC Module (pc.cc)}

The PC module serves dual roles:

\textbf{Source PC (isSource=true):}
\begin{itemize}
    \item Generates periodic traffic (every 2 seconds)
    \item Implements its own Q-learning to select best router
    \item Tracks packet statistics and gate usage
\end{itemize}

\textbf{Destination PC (isDestination=true):}
\begin{itemize}
    \item Receives incoming requests
    \item Immediately sends responses back through arrival gate
    \item Maintains delivery statistics
\end{itemize}

\subsection{Q-Learning Implementation Details}

\subsubsection{Initialization}
All Q-values are initialized to $Q_{\text{init}} = 1.0$ seconds, representing high uncertainty. This encourages exploration of all paths initially.

\subsubsection{Learning Rate}
We use $\alpha = 0.5$, which provides balanced learning:
\begin{itemize}
    \item Fast initial learning from new observations
    \item Stability in steady-state conditions
    \item Each update moves halfway toward the measured value
\end{itemize}

\subsubsection{Packet Tracking}
Each forwarded packet is tracked with:
\begin{itemize}
    \item Unique packet ID
    \item Send timestamp
    \item Destination address
    \item Selected output gate
\end{itemize}

When a response returns, the router:
\begin{enumerate}
    \item Retrieves the original packet information
    \item Calculates RTT
    \item Updates the corresponding Q-value
    \item Removes the tracking entry
\end{enumerate}

\subsection{Implementation Challenges and Solutions}

\subsubsection{Challenge 1: Multiple Gate Support}
\textbf{Problem}: Source PC needed to connect to multiple routers simultaneously.

\textbf{Solution}: Changed gate definition from single \texttt{ppp} to gate array \texttt{ppp[3]}, enabling parallel connections.

\subsubsection{Challenge 2: Response Routing}
\textbf{Problem}: Destination needed to send responses back through the correct router.

\textbf{Solution}: Used \texttt{getArrivalGate()->getIndex()} to determine incoming gate and send response through the same gate.

\subsubsection{Challenge 3: Q-value Convergence}
\textbf{Problem}: Initial implementation showed Q-values not converging.

\textbf{Solution}: Implemented proper packet tracking and RTT measurement instead of relying on Q-UPDATE messages between routers.

\section{Results and Discussion}

\subsection{Experimental Setup}

\begin{itemize}
    \item \textbf{Simulation Duration}: 100 seconds
    \item \textbf{Traffic Pattern}: Periodic packets every 2 seconds (50 packets total)
    \item \textbf{Learning Rate}: $\alpha = 0.5$
    \item \textbf{Exploration Rate}: $\epsilon = 0.1$
    \item \textbf{Initial Q-value}: $Q_{\text{init}} = 1.0$
\end{itemize}

\subsection{Learning Dynamics}

\subsubsection{Q-value Convergence}

The simulation demonstrated clear convergence patterns:

\textbf{Router 1 (FastPath):}
\begin{equation*}
Q: 1.0 \rightarrow 0.51 \rightarrow 0.265 \rightarrow 0.1425 \rightarrow 0.08125 \rightarrow \ldots \rightarrow 0.02
\end{equation*}

\textbf{Router 2 (MediumPath):}
\begin{equation*}
Q: 1.0 \rightarrow 0.55 \rightarrow 0.325 \rightarrow 0.2125 \rightarrow \ldots \rightarrow \sim 0.10
\end{equation*}

\textbf{Router 3 (SlowPath):}
\begin{equation*}
Q: 1.0 \rightarrow 0.6 \rightarrow 0.4 \rightarrow 0.3 \rightarrow \ldots \rightarrow \sim 0.20
\end{equation*}

The convergence follows the expected exponential decay pattern:
\begin{equation}
Q_n = (0.5)^n \cdot Q_0 + [1 - (0.5)^n] \cdot \text{RTT}_{\text{true}}
\end{equation}

\subsubsection{Source PC Learning}

The source PC also learned optimal router selection:

\begin{table}[h]
\centering
\begin{tabular}{|c|c|c|c|}
\hline
\textbf{Gate} & \textbf{Router} & \textbf{Usage} & \textbf{Final Q-value} \\
\hline
0 & Router1 (Fast) & 48/50 (96\%) & 0.022 \\
1 & Router2 (Medium) & 1/50 (2\%) & 0.551 \\
2 & Router3 (Slow) & 1/50 (2\%) & 0.601 \\
\hline
\end{tabular}
\caption{Source PC Gate Usage Statistics}
\label{tab:usage}
\end{table}

\subsection{Performance Metrics}

\begin{itemize}
    \item \textbf{Delivery Success Rate}: 100\% (50/50 packets delivered)
    \item \textbf{Optimal Path Usage}: 96\% (48/50 packets via fastest router)
    \item \textbf{Convergence Time}: $\sim$50 seconds (approximately 25 packets)
    \item \textbf{Exploration Overhead}: 4\% (consistent with $\epsilon = 0.1$ policy)
\end{itemize}

\subsection{Analysis}

\subsubsection{Why 96\% Instead of 90\%?}

The expected exploitation rate with $\epsilon = 0.1$ is 90\%, but we observed 96\%. This discrepancy arises because:

\begin{enumerate}
    \item Initial exploration used systematic selection (3 packets, 1 per gate)
    \item Only 47 remaining packets subject to epsilon-greedy
    \item Random exploration sometimes randomly selects the optimal path
    \item Statistical variance in small sample size (50 packets)
\end{enumerate}

\subsubsection{Convergence Rate}

The learning rate $\alpha = 0.5$ provides rapid initial learning:
\begin{itemize}
    \item After 1 update: Error reduced by 50\%
    \item After 5 updates: Error reduced to $<$5\%
    \item After 10 updates: Error reduced to $<$0.1\%
\end{itemize}

This aggressive learning is suitable for stable networks but might be too reactive for highly dynamic environments.

\subsubsection{Exploration-Exploitation Trade-off}

The epsilon-greedy policy with $\epsilon = 0.1$ balances:
\begin{itemize}
    \item \textbf{Exploitation}: Uses best known path 90\% of time (efficiency)
    \item \textbf{Exploration}: Tests alternatives 10\% of time (adaptability)
\end{itemize}

This continued exploration enables the system to detect if network conditions change.

\subsection{Validation Against Theory}

Our results validate theoretical predictions:

\begin{enumerate}
    \item \textbf{Bellman Optimality}: Q-values converge to true delivery times
    \item \textbf{Policy Convergence}: System converges to near-optimal policy (96\% optimal)
    \item \textbf{Exploration Necessity}: Random exploration maintains system adaptability
    \item \textbf{Learning Speed}: Convergence achieved within reasonable time frame
\end{enumerate}

\section{Figures}

\begin{figure}[h]
\centering
% Placeholder for Q-value convergence graph
\fbox{\parbox{0.8\textwidth}{\centering\vspace{3cm}[Q-value Convergence Graph - To be added]\vspace{3cm}}}
\caption{Q-value convergence over time for all three routers showing learning dynamics}
\label{fig:convergence}
\end{figure}

\section{Conclusion}

This project successfully implemented and validated Q-Routing, a reinforcement learning-based adaptive routing protocol, in OMNeT++. The implementation demonstrates several key findings:

\subsection{Key Achievements}

\begin{enumerate}
    \item \textbf{Successful Learning}: Routers autonomously learned optimal routing decisions without prior knowledge of network topology or link characteristics.
    
    \item \textbf{Rapid Convergence}: Q-values converged to reflect true end-to-end delays within approximately 50 seconds (25 packets) of simulation time.
    
    \item \textbf{Optimal Path Selection}: The system achieved 96\% utilization of the fastest available path, exceeding the theoretical 90\% expected from epsilon-greedy policy.
    
    \item \textbf{100\% Reliability}: All packets were successfully delivered, demonstrating system robustness.
    
    \item \textbf{Adaptive Capability}: Continued 10\% exploration maintains the ability to detect and adapt to changing network conditions.
\end{enumerate}

\subsection{Insights}

\begin{itemize}
    \item \textbf{Smart Exploration}: The two-phase exploration strategy (systematic then epsilon-greedy) ensures thorough initial exploration while maintaining long-term adaptability.
    
    \item \textbf{Learning Rate Impact}: The choice of $\alpha = 0.5$ provided good balance between fast learning and stability.
    
    \item \textbf{Simplified Q-Routing}: For single-hop topologies, using direct RTT measurement (Equation \ref{eq:simple}) is more efficient than full Q-value propagation.
\end{itemize}

\subsection{Practical Applications}

Q-Routing principles can be applied to:
\begin{itemize}
    \item Software-Defined Networks (SDN) for dynamic flow routing
    \item Content Delivery Networks (CDN) for optimal server selection
    \item Mobile ad-hoc networks (MANET) with frequently changing topology
    \item Data center networks with varying load conditions
    \item IoT networks with intermittent connectivity
\end{itemize}

\subsection{Limitations and Future Work}

\textbf{Current Limitations:}
\begin{itemize}
    \item Single source-destination pair
    \item Static link delays (no dynamic changes during simulation)
    \item Simple topology (parallel paths only)
    \item No congestion modeling
\end{itemize}

\textbf{Future Extensions:}
\begin{enumerate}
    \item \textbf{Dynamic Network Conditions}: Implement time-varying link delays and failures to test adaptability
    
    \item \textbf{Multi-hop Topology}: Extend to complex topologies requiring full Q-value propagation between routers
    
    \item \textbf{Traffic Models}: Introduce realistic traffic patterns with varying loads and priorities
    
    \item \textbf{Scalability Testing}: Evaluate performance with larger topologies (10+ routers, multiple sources)
    
    \item \textbf{Comparative Analysis}: Compare Q-Routing performance against OSPF, EIGRP, and other protocols
    
    \item \textbf{Advanced RL Algorithms}: Implement Deep Q-Networks (DQN) or other modern RL techniques
    
    \item \textbf{Multi-objective Optimization}: Extend to consider multiple metrics (delay, bandwidth, energy)
\end{enumerate}

\subsection{Final Remarks}

This project demonstrates that reinforcement learning can effectively solve the routing problem in computer networks. Q-Routing's ability to learn from experience without requiring global knowledge or manual configuration makes it particularly attractive for modern dynamic network environments. The successful implementation in OMNeT++ provides a foundation for further research into learning-based network protocols and serves as an educational tool for understanding the intersection of machine learning and networking.

The high performance achieved (96\% optimal path usage, 100\% delivery rate) validates the theoretical foundations of Q-learning applied to routing and suggests promising directions for practical deployment in adaptive network systems.

\begin{thebibliography}{9}

\bibitem{boyan1994packet}
J. A. Boyan and M. L. Littman, 
``Packet routing in dynamically changing networks: A reinforcement learning approach,''
in \textit{Advances in Neural Information Processing Systems 6 (NIPS 1993)}, 
J. D. Cowan, G. Tesauro, and J. Alspector, Eds. 
Morgan Kaufmann Publishers, 1994, pp. 671--678.

\bibitem{sutton2018reinforcement}
R. S. Sutton and A. G. Barto,
\textit{Reinforcement Learning: An Introduction}, 2nd ed.
Cambridge, MA: MIT Press, 2018.

\bibitem{watkins1992q}
C. J. C. H. Watkins and P. Dayan,
``Q-learning,''
\textit{Machine Learning}, vol. 8, no. 3, pp. 279--292, May 1992.

\bibitem{kumar2004comparative}
S. Kumar and R. Miikkulainen,
``Dual reinforcement Q-routing: An on-line adaptive routing algorithm,''
in \textit{Proc. Artificial Neural Networks in Engineering (ANNIE 2004)},
St. Louis, MO, USA, Nov. 2004.

\bibitem{choi2001q}
S. P. M. Choi and D. Y. Yeung,
``Predictive Q-routing: A memory-based reinforcement learning approach to adaptive traffic control,''
in \textit{Advances in Neural Information Processing Systems 8 (NIPS 1995)},
D. S. Touretzky, M. C. Mozer, and M. E. Hasselmo, Eds.
Cambridge, MA: MIT Press, 1996, pp. 945--951.

\bibitem{varga2010omnet}
A. Varga,
``OMNeT++,''
in \textit{Modeling and Tools for Network Simulation},
K. Wehrle, M. G{\"u}ne{\c{s}}, and J. Gross, Eds.
Berlin, Germany: Springer, 2010, pp. 35--59.

\bibitem{mellouk2009qos}
A. Mellouk, Ed.,
\textit{Quality of Service Mechanisms in Next Generation Heterogeneous Networks}.
London, UK: ISTE-Wiley, 2009.

\bibitem{littman2015reinforcement}
M. L. Littman,
``Reinforcement learning improves behaviour from evaluative feedback,''
\textit{Nature}, vol. 521, no. 7553, pp. 445--451, May 2015.

\end{thebibliography}

\end{document}
